% Standalone TikZ figure: a small feedforward neural network built out of
% the perceptron unit of Figure~\ref{fig:perceptron}.
% \input from the chapter's raw-LaTeX block; not rendered by Obsidian.
\begin{figure}[H]
\centering
\begin{tikzpicture}[
    neuron/.style={circle, draw, minimum size=8mm},
    inlayer/.style={neuron, fill=mondrianblue!15},
    hidlayer/.style={neuron, fill=mondrianyellow!25},
    outlayer/.style={neuron, fill=mondrianred!20},
    >=stealth
]
    \foreach \i in {1,...,4}
        \node[inlayer] (I-\i) at (0, {-\i*1.1+2.75}) {};
    \foreach \i in {1,...,5}
        \node[hidlayer] (H-\i) at (3.2, {-\i*1.1+3.85}) {};
    \foreach \i in {1,...,2}
        \node[outlayer] (O-\i) at (6.4, {-\i*1.1+2.2}) {};

    \foreach \i in {1,...,4}
        \foreach \j in {1,...,5}
            \draw[->, gray, thin] (I-\i) -- (H-\j);
    \foreach \i in {1,...,5}
        \foreach \j in {1,...,2}
            \draw[->, gray, thin] (H-\i) -- (O-\j);

    \node at (0,3.35) {\small input layer};
    \node at (3.2,4.45) {\small hidden layer};
    \node at (6.4,2.8) {\small output layer};
\end{tikzpicture}
\caption{A feedforward neural network (multilayer perceptron). Each node is a perceptron unit as in Figure~\ref{fig:perceptron}, organized into layers, with every neuron of one layer feeding into every neuron of the next.}
\label{fig:neural-network}
\end{figure}
